\chapter{Introduction}\label{ch:introduction}

\app\ is an interactive single-line diagram editor with a power system analysis engine
behind it. You draw a network on a canvas, much as in DIgSILENT PowerFactory, PSS/E or
ETAP. \app\ converts it to per unit and runs the studies a power system analysis course
asks for, with the results printed the way the course prints them. It started from the
KTH course \emph{EG2100 Power System Analysis}. Version~1.0 solves all four course
assignments (S1, S2, S3 and LF) and reproduces their reference solutions digit by digit.

\section{What \app\ does}

\begin{description}
  \item[Draw] buses, external grids, generators, two-winding transformers, PI-model lines
    and cables, loads, shunt capacitors and reactors, static var compensators and raw
    impedances. Wires route themselves at right angles. Busbars resize, labels stay
    upright, and every edit can be undone (\cref{ch:drawing,ch:components}).
  \item[Parametrise] any numeric field with an expression over project variables, such as
    \code{(217 + A)/1000}. Change \code{A} and the whole network is re-evaluated
    (\cref{ch:parameters}).
  \item[Organise] the work in \emph{study cases}. A case is a study, its options and,
    optionally, a variant of the network: elements taken out of service and fields
    overridden by an expression. The drawn network is never changed (\cref{ch:cases}).
  \item[Analyse] with eleven studies (\cref{ch:studies}):
    \begin{itemize}
      \item the per-unit model and the sequence model;
      \item two-port (ABCD) matrices, equivalent impedances and linear bus voltages;
      \item the positive-, negative- and zero-sequence admittance matrices;
      \item Thevenin equivalents in all three sequences;
      \item three-phase, single line-to-ground, line-to-line and double line-to-ground
        faults through fault and ground impedances;
      \item unbalanced loads;
      \item the Newton--Raphson load flow, with its iteration log and Jacobian.
    \end{itemize}
  \item[Read] each result in its own tab: a headline, warnings and tables formatted like
    MATLAB's \code{format short}. Annotate the diagram with any result and colour it by
    voltage level or bus voltage (\cref{ch:results}).
  \item[Edit in bulk] every element of a type in the data manager, a spreadsheet-like
    table (\cref{ch:data-manager}).
  \item[Export] results as text, CSV, LaTeX tables or JSON, the diagram as PNG, SVG or
    PDF, and matrices as MATLAB literals that keep every digit (\cref{ch:export}).
  \item[Automate] from a terminal with \file{run\_headless.py}, which runs a project's
    cases, sets variables and exports what it ran (\cref{ch:console}).
\end{description}

\section{Principles}

Four ideas run through the program, and they explain most of its behaviour.

\begin{description}
  \item[Glass box.] Every number can be traced. The per-unit table shows each component's
    conversion with its base factor. The admittance, impedance and Jacobian matrices are
    printed with row and column labels, and the load flow logs every iteration.
  \item[Verified.] Each study checks itself by a second, independent method. The Thevenin
    impedance is computed from the impedance matrix and again by Kron reduction. The
    two-port matrix comes from the element cascade and from the Y-bus. The load flow
    closes the power balance. The cross-check residuals are part of the results.
  \item[Parametric.] Variables, expressions and study cases let one project hold every
    variant of a problem, and re-evaluate it when a parameter changes.
  \item[Course-grade presentation.] Impedance and complex power are printed in rectangular
    form, voltage and current in polar form with the angle in degrees, and matrices with
    MATLAB's common factor (\code{1.0e+02 *}).
\end{description}

\section{About this manual}

\Cref{ch:getting-started} gets you from installation to a first solved network in a few
minutes. \Cref{ch:window,ch:drawing} explain the window and the editor, and
\cref{ch:components} is the reference for every component and field. Then come per-unit
bases (\cref{ch:bases}), parametric variables (\cref{ch:parameters}) and study cases
(\cref{ch:cases}). \Cref{ch:studies} explains each study: what it computes, its options
and its output. \Cref{ch:results,ch:data-manager,ch:export,ch:console} cover reading,
editing in bulk, exporting and the command line. \Cref{ch:assignments} walks through the
four EG2100 assignments task by task, \cref{ch:files} describes the project file, and
\cref{ch:troubleshooting} lists the messages you may meet and what to do about them. The
appendices collect the model equations, the keyboard shortcuts and how the program is
verified.

Conventions used in this manual:

\begin{itemize}
  \item Labels on the screen are set like \gui{Properties}. Menu paths look like
    \menu{File > Export results...}, and keys like \keys{\ctrl + S}.
  \item Field names, option names and values are set like \code{p\_mw} or
    \code{constant\_power}. These are the names that bindings, overrides and the project
    file use.
  \item File names look like \file{examples/lf\_220kv.psn}, and commands typed in a
    terminal like \cli{python run\_gui.py}.
  \item The reference tables of \cref{ch:components,ch:studies} are generated from the
    program's own field definitions, so they always match the version this manual
    describes.
\end{itemize}

\section{Scope of version 1.0}

\app~1.0 models the balanced three-phase network in per unit, with symmetrical components
for unbalanced conditions. It does not yet model breakers and switches as components,
three-winding transformers, off-nominal taps or phase shifters, transformer phase shifts
in the sequence networks, IEC~60909 short circuits (voltage factors), a distributed slack,
or dynamics. \Cref{sec:limitations} lists the limitations, and \file{docs/ROADMAP.md} in
the repository lists what comes next.
